% Instantiation Operator
\section{Instantiation Operator}

\subsection{Definition}

The instantiation operator $\text{INST}(\sigma)$ introduces an entity or object into the operational context. This is the fundamental operation that brings a referent into the scope of subsequent operations.

\subsection{Formal Characteristics}

\begin{itemize}
    \item \textbf{Position}: Predominantly initial or early in sequence
    \item \textbf{Cardinality}: Typically appears once per inscription
    \item \textbf{Independence}: Can stand alone (complete inscription)
    \item \textbf{Productivity}: High frequency of use
\end{itemize}

\subsection{Sign Classes Functioning as Instantiation}

\subsubsection{Primary Instantiation Signs}
\begin{itemize}
    \item Fish signs (M361--M372): Domain or institutional identifiers
    \item Jar/vessel signs (M346--M360): Object type identifiers
    \item Animal signs (M391--M410): Commodity or domain classifiers
\end{itemize}

\subsubsection{Secondary Instantiation}
Some signs can function as instantiation in specific contexts:
\begin{itemize}
    \item Anthropomorphic signs (role/office instantiation)
    \item Tool signs (operational domain instantiation)
\end{itemize}

\subsection{Operational Semantics}

The instantiation operator establishes:
\begin{enumerate}
    \item \textbf{Type}: What category of entity is being referenced
    \item \textbf{Scope}: The domain for subsequent operations
    \item \textbf{Context}: Administrative, commercial, or institutional frame
\end{enumerate}

\subsection{Examples from Corpus}

Typical patterns:
\begin{itemize}
    \item [FISH] = Complete inscription (institutional marker alone)
    \item [FISH][STROKES] = Instantiation + counting
    \item [JAR] = Object instantiation
    \item [JAR][STROKES][TERMINAL] = Object + quantity + closure
\end{itemize}

\subsection{Theoretical Interpretation}

In computational terms, instantiation corresponds to:
\begin{itemize}
    \item Object creation in programming
    \item Variable declaration
    \item Type specification
    \item Context initialization
\end{itemize}

This operation is \emph{declarative}: it states the existence and type of an entity without necessarily performing actions upon it.
